%!TEX root = ../../main.tex
\section{승률과 시간의 기준선}
\label{sec:pred-baselines}

M-RQ3의 핵심은 예측이 가능한가가 아니라, 교전 전에 누구나 아는 두 양, 즉 현재 승률 $\ppre$와 경기 시각 $t$를 \emph{넘어서} 예측 가능한가이다. 따라서 기준선은 이 두 양의 함수여야 하고, 충분히 유연해야 한다.

\begin{itemize}[leftmargin=1.5em]
  \item $\PTlin$: $\sigma(a + b\,\ppre + c\,t)$, 표준화된 $[\ppre, t]$ 위의 로지스틱 회귀.
  \item $\PTflex$: $\ppre$와 분 단위 시각 $t$ 각각에 대한 3차 B-스플라인 기저(균일 매듭, 상수 외삽; 매듭 수 $\knotsP$, $\knotsT$)와 두 기저의 완전 텐서곱을 표준화한 뒤 $L_2$ 로지스틱 회귀(lbfgs)로 가중치 $1/n_m$ 아래 적합한 모델. 탐색 격자는 $\ppre$ 매듭 수 $\{4, 6\} \times C \in \{0.01, 0.1, 1\}$이었고 Q\_SELECT의 경기 가중 Brier가 고른 승자는 $(\knotsP, \knotsT, \ridgeC)$이다. 이 설정은 격자에서 가장 강하게 정규화된 모서리이다.
  \item $b(p)$: $\ppre$만의 스플라인 기준선(매듭 \knotsPonly, $C = \ridgeC$). 진단용이다.
\end{itemize}

실행이 서면 설계와 달랐던 부분은 다음과 같고 역사를 고쳐 쓰지 않는다. 매듭은 가중 분위수가 아니라 균일하고, 적합 가중치는 평균 정규화되지 않았으며, 하이퍼파라미터가 항등/시그모이드 보정기보다 먼저 선택되었다(두 단계). S에 대해서는 같은 격자를 S Q\_SELECT에서 탐색하여 같은 승자 $(\knotsP, \knotsT, \ridgeC)$가 선정되었고, $\PTlinS$와 $b(p)_S$는 S TRAIN에서 같은 방식으로 적합되었다. RR12의 두 단계 규칙을 S에 적용하면 모든 모델($\qS$, $\PTflexS$, $\PTlinS$, $b(p)_S$)에 양의 기울기 시그모이드가 선택되었으나, 규모분리 계약이 항등 보정기를 사전에 고정하였고 $\qTtoS$를 동결 T 모델의 그대로 적용으로 정의하였으므로, 본 논문의 모든 S 결과는 \emph{항등} 보정기를 쓰고 선택 변형은 민감도 행으로만 보고한다.
